\section{Introdução}

A otimização de inferência de modelos de aprendizado de máquina representa
uma interseção entre os algoritmos de ciência de dados e as exigências
práticas de engenharia de sistemas \cite{nakandala2020}. Enquanto a
comunidade acadêmica dedica grande atenção ao treinamento de modelos, a
entrega eficiente de predições em produção permanece subexplorada,
especialmente para \textit{ensembles} de árvores de decisão
\cite{asadi2024}.

Existe um abismo entre treinar um modelo em um \textit{Jupyter Notebook} e
colocá-lo em produção com latência previsível e escalabilidade. Em domínios
que servem dados tabulares de forma síncrona, como detecção de fraude em
tempo real, escoragem de crédito e \textit{Real-Time Bidding} (RTB) em
publicidade digital, cada requisição exige uma decisão em milissegundos para
bloquear ou aprovar uma operação, depende de uma resposta instantânea de
aprovação ou impõe orçamentos de latência inferiores a 10\,ms. Nessas
condições, a latência de cauda no percentil P99 não é uma métrica acadêmica
e sim a diferença entre capturar ou perder uma transação.

Modelos de \textit{gradient boosting} como o XGBoost dominam essas
aplicações de dados tabulares, superando redes neurais em benchmarks
representativos \cite{chen2016, grinsztajn2022}. Entretanto, servir esses
modelos com garantias de latência é um desafio técnico distinto: a
travessia condicional de centenas de árvores binárias causa padrões
irregulares de acesso à memória \cite{nakandala2020}, e o sobrecusto da
camada de \textit{serving} ao redor da inferência tende a dominar o tempo de
resposta para modelos cuja inferência é da ordem de microssegundos.

Em arquiteturas modernas de microsserviços, sistemas de aprendizado de
máquina são frequentemente implantados em contêineres orquestrados via
Docker ou Kubernetes \cite{zaharia2024}. Nesses ambientes, \textit{cold
starts} (inicialização de novas instâncias) e \textit{jitter} de sistema,
provocado por pausas do \textit{garbage collector}, trocas de contexto de
CPU e \textit{cache misses}, tornam a mediana da latência (P50) uma métrica
inadequada para avaliar a confiabilidade operacional \cite{dean2013}. A
latência de cauda, representada pelo percentil P99, captura as anomalias
severas que afetam a experiência dos usuários, e medi-la corretamente exige
um gerador de carga de modelo aberto que não desacelere quando o servidor
engasga \cite{tene2015}.

Este trabalho conduz uma análise quantitativa comparativa entre dois
pipelines de entrega para um modelo XGBoost treinado sobre o conjunto de
dados HIGGS \cite{higgs2014} e exportado em formato ONNX
\cite{onnxruntime2024}. O primeiro pipeline emprega FastAPI com ONNX
Runtime, representando a abordagem de servidor web generalista; o segundo
emprega BentoML com \textit{adaptive batching}, representando a abordagem
especializada que agrupa requisições concorrentes em lotes antes da
inferência. Ambos os pipelines são conteinerizados, sujeitos a limites
idênticos de CPU e memória, e submetidos a uma varredura matricial sobre os
eixos de paralelismo do servidor (\textit{workers} e \textit{threads}) e
intensidade de carga (RPS).

As contribuições deste trabalho são as seguintes. Primeiro, uma comparação
empírica em 219 células experimentais cobrindo quatro variantes de
configuração e cinco taxas de requisições por segundo, com três repetições
por célula. Segundo, uma decomposição da latência ponta a ponta em três
camadas de medição (inferência pura, HTTP no host e contêiner Docker
completo), permitindo atribuir o sobrecusto a componentes específicos da
plataforma de \textit{serving}. Terceiro, uma métrica de vazão útil sob SLA que
permite comparar de forma justa um sistema que descarta requisições por
prazo expirado contra um sistema que as enfileira como latência adicional.
Quarto, evidência direta de utilização de CPU coletada em um subconjunto de
mecanismo de 60 execuções, identificando a corrotina de \textit{dispatcher}
do BentoML como o ponto de serialização que limita sua vazão.

Os resultados refutam a hipótese inicial de que o \textit{adaptive batching}
absorveria concorrência com mais elegância e comprimiria a latência de
cauda. O FastAPI sustentou entre 99,7\% e 100\% de vazão útil dentro de um
SLA de 50\,ms sobre toda a matriz de 50 a 1000 RPS, enquanto o BentoML
colapsou entre 350 e 600 RPS escalando com o número de \textit{workers}, e o
P99 cresceu duas ordens de magnitude sobre o mesmo intervalo de RPS. A
evidência de CPU mostra que o BentoML consome aproximadamente 2,5 vezes a
CPU do FastAPI sob carga idêntica, e mesmo assim colapsa: o gargalo é a
serialização imposta pela única corrotina de \textit{dispatcher} por
\textit{worker}, não escassez de recursos computacionais. Para inferência
barata de árvores em CPU sobre uma linha por requisição, o sobrecusto por
requisição do \textit{adaptive batching} excede qualquer amortização de
lote possível para um custo de inferência da ordem de 14 microssegundos.

O restante desta monografia está organizado em oito seções. A Seção~2
revisa o operador ONNX para \textit{ensembles} de árvores, a teoria de
latência de cauda e \textit{Coordinated Omission}, e o desenho do
\textit{adaptive batching}. A Seção~3 posiciona o trabalho frente a
benchmarks publicados de inferência de florestas de decisão e comparações
de plataformas de \textit{serving}. A Seção~4 descreve a metodologia
experimental: decomposição em três camadas, varredura matricial, geração
de carga com Vegeta, pinagem de CPU e procedimentos estatísticos. A
Seção~5 documenta a implementação dos dois servidores e da infraestrutura
de \textit{benchmarking}. A Seção~6 apresenta os resultados quantitativos. A
Seção~7 discute o mecanismo identificado, os regimes em que o
\textit{adaptive batching} ainda seria vantajoso e as ameaças à validade. A
Seção~8 conclui e aponta trabalhos futuros.

